\documentclass[conference]{IEEEtran}

\IEEEoverridecommandlockouts
% The preceding line is only needed to identify funding in the first footnote. If that is unneeded, please comment it out.
\usepackage{cite}
\usepackage{amsmath,amssymb,amsfonts}
\usepackage{algorithmic}
\usepackage{graphicx}
\usepackage{textcomp}
\usepackage{xcolor}
\bibliographystyle{IEEEtran}

\def\BibTeX{{\rm B\kern-.05em{\sc i\kern-.025em b}\kern-.08em
    T\kern-.1667em\lower.7ex\hbox{E}\kern-.125emX}}
\begin{document}

\title{Traffic Signal Violation Detection at Simpang Gedongan Sleman Using HOG-SVM Classifier, HSV-Based Segmentation, and SORT}

\author{\IEEEauthorblockN{1\textsuperscript{st} Deira Aisya Refani}
\text{24/532821/PA/22539}\\
\IEEEauthorblockA{\textit{Universitas Gadjah Mada} \\
\textit{Department of Computer Science and Electronics}\\
Yogyakarta, Indonesia \\
deiraaisyarefani532821@mail.ugm.ac.id}
\and
\IEEEauthorblockN{2\textsuperscript{nd} Finanazwa Ayesha}
\text{24/532953/PA/22556}\\
\IEEEauthorblockA{\textit{Universitas Gadjah Mada} \\
\textit{Department of Computer Science and Electronics}\\
Yogyakarta, Indonesia \\
email}
\and
\IEEEauthorblockN{3\textsuperscript{rd} Farrel Tsaqif Anindyo}
\text{24/536735/PA/22769}\\
\IEEEauthorblockA{\textit{Universitas Gadjah Mada} \\
\textit{Department of Computer Science and Electronics}\\
Yogyakarta, Indonesia \\
email}
\and
\IEEEauthorblockN{4\textsuperscript{th} Yohana Butar Butar}
\text{24/546690/PA/23212}\\
\IEEEauthorblockA{\textit{Universitas Gadjah Mada} \\
\textit{Department of Computer Science and Electronics}\\
Yogyakarta, Indonesia \\
email}
}

\maketitle


\section{Introduction}
This document is a model and instructions for \LaTeX.
Please observe the conference page limits. 

\subsection{Background}
blabla
\subsection{Problem Statement}
blabka
\subsection{Objective}
blabla
\subsection{Scope}
This study is limited to traffic signal violation detection using static CCTV footage at Simpang Gedongan Sleman, Yogyakarta. The system utilizes HOG-SVM for vehicle detection and HSV-based segmentation for traffic light detection. The analysis is restricted to a fixed camera view and predefined regions of interest, without considering varying environmental conditions or multiple camera setups.

\section{Literature Review}
\subsection{Related Work}
\subsection{HSV (Hue, Saturation, Value) Color Model}
\subsection{Histogram of Oriented Gradients (HOG)}
\subsection{Support Vector Machine (SVM)}
\subsection{Simple Online and Realtime Tracking (SORT)}

\section{Methodology}
\subsection{Proposed Method} 
\subsection{Dataset}
\subsection{Preprocessing}
\subsection{Vehicle Detection using HOG-SVM}
\subsection{Vehicle Tracking using SORT}
\subsection{Traffic Light Detection using HSV}
\subsection{Violation Detection}
\subsection{Evaluation Metrics}
\subsubsection{Accuracy}
\subsubsection{Precision}
\subsubsection{Recall}
\subsubsection{F1-Score}

\bibliography{references}  

\end{document}